\section*{Capítulo \ref{ch:b2:atomic-orbitals} --- Orbitales atómicos: formas, energías y tamaños}

\begin{solution}{exo:b2:atomic-orbitals:1}
3p: $n = 3$, $l = 1$; un \omterm{def:b2:atomic-orbitals:node}{nodo radial} ($n - l - 1$) y un \omterm{def:b2:atomic-orbitals:node}{nodo angular}. 4d: $n = 4$,
$l = 2$; un \omterm{def:b2:atomic-orbitals:node}{nodo radial} y dos \omterm{def:b2:atomic-orbitals:node}{nodos angulares}.
\end{solution}

\begin{solution}{exo:b2:atomic-orbitals:2}
$P_{2p} \propto r^4\mathrm e^{-r}$ ($r$ en $a_0$): $d\ln P/dr = 4/r - 1 = 0$ para $r = 4$:
$4a_0 = \qty{212}{pm}$.
\end{solution}

\begin{solution}{exo:b2:atomic-orbitals:3}
$1 - \mathrm e^{-2}(1 + 2 + 2) = 1 - 5\mathrm e^{-2} = 0.323$.
\end{solution}

\begin{solution}{exo:b2:atomic-orbitals:4}
Cuatro lóbulos que apuntan entre los ejes $x$ e $y$, en el plano $xy$, de signos
alternos (positivos donde $xy > 0$); los planos nodales son $xz$ e $yz$. \omterm{def:b2:atomic-orbitals:node}{Nodos radiales}:
$n - l - 1 = 0$.
\end{solution}

\begin{solution}{exo:b2:atomic-orbitals:5}
Oxígeno, $(1\mathrm s)^2(2\mathrm s, 2\mathrm p)^6$: $\sigma = 5 \times 0.35 + 2 \times 0.85
= 3.45$, $Z^* = 4.55$; radio $4/4.55 = 0.88a_0$, \qty{47}{pm}.
\end{solution}

\begin{solution}{exo:b2:atomic-orbitals:6}
La ionización arranca los electrones menos ligados. Los electrones 4s
(\qty{-14}{eV}) están mucho menos ligados que los 3d (\qty{-59}{eV}): se van primero los dos
electrones 4s y luego un electrón 3d, lo que da $[\ce{Ar}]\,3\mathrm d^5$.
\end{solution}

\begin{solution}{exo:b2:atomic-orbitals:7}
Litio, 2s: $Z^* = 3 - 2 \times 0.85 = 1.30$, $I = 13.6 \times 1.30^2/4 =
\qty{5.75}{eV}$ (medida, 5.39: un 7~\% demasiado alta). Sodio, 3s: $Z^* = 11 - 8 \times 0.85 -
2 \times 1.00 = 2.20$, $I = 13.6 \times 2.20^2/9 = \qty{7.3}{eV}$ (medida, 5.14): las
reglas, un modelo aproximado, sobrestiman la energía de enlace del electrón externo, cada vez
más al bajar por el grupo.
\end{solution}

\begin{solution}{exo:b2:atomic-orbitals:8}
Sodio: $Z^* = 2.20$, radio $9/2.20 = 4.1a_0$ (\qty{216}{pm}). Cloro, 3p:
$\sigma = 6 \times 0.35 + 8 \times 0.85 + 2 \times 1.00 = 10.90$, $Z^* = 6.10$, radio
$9/6.10 = 1.5a_0$ (\qty{78}{pm}). La misma capa, pero el electrón externo del cloro
siente casi el triple de carga, y es atraído hacia dentro.
\end{solution}

\begin{solution}{exo:b2:atomic-orbitals:9}
Ambos tienen la \omterm{def:b2:atomic-orbitals:radial-angular}{parte angular} $1/\sqrt{4\pi}$, de modo que la integral se reduce a
\[
  \int_0^\infty R_{1s}R_{2s}r^2\,dr \propto \int_0^\infty (2 - r)r^2\mathrm e^{-3r/2}\,dr =
  2 \times \frac{2!}{(3/2)^3} - \frac{3!}{(3/2)^4} = \frac{32}{27} - \frac{32}{27} = 0 .
\]
\end{solution}

\begin{solution}{exo:b2:atomic-orbitals:10}
$\int_0^\infty N^2r^2\mathrm e^{-r}r^2\,dr = N^2 \times 4! = 1$: $N = 1/\sqrt{24} =
1/(2\sqrt6)$, como en la tabla.
\end{solution}

\begin{solution}{exo:b2:atomic-orbitals:11}
$\langle r\rangle = 1.5a_0$, \omterm{def:b2:atomic-orbitals:radial-distribution}{radio más probable} $a_0$. $P(r) = 4r^2\mathrm e^{-2r}$ sube
bruscamente desde cero y baja lentamente: la larga cola a $r$ grande pesa más en la
media que en la posición del máximo.
\end{solution}

\begin{solution}{exo:b2:atomic-orbitals:12}
$Y_{p_z} \propto z/r = \cos\theta$ se anula para $\theta = \pi/2$: el plano $xy$.
$\psi$ tiene el signo de $z$, positivo por encima y negativo por debajo. Sobre el plano la
\omterm{def:b2:atomic-orbitals:wavefunction}{densidad de probabilidad} es nula.
\end{solution}

\begin{solution}{pb:b2:atomic-orbitals:1}
\textbf{1.} $\int|\psi|^2\,dV = \frac{4\pi}{\pi a_0^3}\int_0^\infty r^2\mathrm e^{-2r/a_0}\,dr =
\frac{4}{a_0^3} \times \frac{a_0^3}{4} = 1$.
\textbf{2.} $P(r) = 4\pi r^2|\psi|^2 = (4/a_0^3)\,r^2\mathrm e^{-2r/a_0}$.
\textbf{3.} $dP/dr = 0$: $2/r - 2/a_0 = 0$, $r = a_0$.
\textbf{4.} $\langle r\rangle = (4/a_0^3)\int_0^\infty r^3\mathrm e^{-2r/a_0}\,dr = (4/a_0^3)
\times 6(a_0/2)^4 = 1.5a_0$.
\textbf{5.} La distribución es asimétrica, con una larga cola a $r$ grande.
\textbf{6.} $|\psi|^2$ es una densidad por unidad de volumen, máxima en el núcleo; $P(r)$
cuenta el volumen de la capa, $4\pi r^2\,dr$, que se anula en $r = 0$.
\textbf{7.} \qty{52.9}{pm}.
\textbf{8.} Con $x = r/a_0$: $\int_{x_0}^\infty 4x^2\mathrm e^{-2x}\,dx$; integrando dos veces por partes,
$= \mathrm e^{-2x_0}(2x_0^2 + 2x_0 + 1)$.
\textbf{9.} $5\mathrm e^{-2} = 0.677$.
\textbf{10.} $13\mathrm e^{-4} = 0.238$.
\textbf{11.} $\mathrm e^{-2x}(1 + 2x + 2x^2) = 0.10$ para $x \approx 2.66$: \qty{141}{pm}.
\textbf{12.} La densidad no se anula a ninguna distancia: no hay borde. Pero el 90~\% de la
probabilidad está dentro de \qty{141}{pm}: un tamaño práctico.
\textbf{13.} 2s: un \omterm{def:b2:atomic-orbitals:node}{nodo radial} (esfera de radio $2a_0$), ningún \omterm{def:b2:atomic-orbitals:node}{nodo angular}. 2p: un
\omterm{def:b2:atomic-orbitals:node}{nodo angular} (un plano), ningún \omterm{def:b2:atomic-orbitals:node}{nodo radial}.
\textbf{14.} En $r = 2a_0$.
\textbf{15.} $4a_0$ (ejercicio 2).
\textbf{16.} $d\ln P/dr = 2/r - 2/(2 - r) - 1 = 0$ da $r^2 - 6r + 4 = 0$, $r = 3 \pm \sqrt5$:
$0.76a_0$ (pequeño máximo interior) y $5.24a_0$ (máximo principal).
\textbf{17.} El 2s, por su máximo interior. En un átomo polielectrónico un electrón 2s
penetra en la capa 1s, está menos apantallado que un electrón 2p y tiene una energía
más baja.
\textbf{18.} Todos los radios se dividen por $Z = 6$: 2p a $0.67a_0$ (\qty{35}{pm}); máximos de 2s
a $0.13a_0$ y $0.87a_0$.
\textbf{19.} C 3.25, N 3.90, O 4.55.
\textbf{20.} $\varepsilon = -13.6\,Z^{*2}/4$: C \qty{-35.9}{eV}, N \qty{-51.7}{eV}, O
\qty{-70.4}{eV}.
\textbf{21.} Como en el capítulo: \qty{11.5}{eV} (medida, 11.26).
\textbf{22.} N: átomo, $5 \times (-13.6 \times 3.90^2/4) = \qty{-258.6}{eV}$; \ce{N+}, $Z^* =
4.25$, $4 \times (-13.6 \times 4.25^2/4) = \qty{-245.7}{eV}$; $I = \qty{12.9}{eV}$.
\textbf{23.} O: átomo, $6 \times (-13.6 \times 4.55^2/4) = \qty{-422.3}{eV}$; \ce{O+}, $Z^* =
4.90$, $5 \times (-13.6 \times 4.90^2/4) = \qty{-408.2}{eV}$; $I = \qty{14.2}{eV}$. Modelo:
11.5, 12.9, 14.2; medidas: 11.26, 14.53, 13.62.
\textbf{24.} La energía medida del oxígeno es menor que la del nitrógeno: en el
oxígeno el cuarto electrón 2p debe compartir un orbital con otro, y su
repulsión facilita arrancarlo, mientras que los tres electrones p desapareados del nitrógeno
forman una subcapa semillena estable. Las \omterm{def:b2:atomic-orbitals:slater}{reglas de Slater} tratan igual a todos los electrones de un grupo
y no saben nada del apareamiento.
\textbf{25.} El electrón 1s del hidrógeno está más lejos que $2a_0$ con una probabilidad
$\boldsymbol{13\mathrm e^{-4} \approx 0.238}$.
\end{solution}
